% 量子力学1 量子力学1-12 問題2（048-free。当時は 量子力学1-6 問題3） 解いた日 2026-11-01
%% kotae
\toi{1}{◎}{}
\begin{align*}
\comm{V}{\hat{p}}\psi&=\bigl(V(x)\hat{p}-\hat{p}V(x)\bigr)\psi\\
&=-\ii\hbar V\psi'+\ii\hbar(V'\psi+V\psi')\\
&=\ii\hbar V'(x)\psi
\end{align*}
\begin{align*}
\comm{\hat{H}}{\hat{p}}&=\frac{1}{2m}\comm{\hat{p}^{2}}{\hat{p}}+\comm{V}{\hat{p}}\\
&=\ii\hbar V'(x)
\end{align*}
\toi{2}{◎}{}
$V=0$ なら $\comm{\hat{H}}{\hat{p}}=0$ より同時に観測可能。
\[ \psi=C\ee^{\ii kx} \]
\[ \hat{p}\psi=\kakimachigai{C\hbar k\psi}{\hbar k\psi},\qquad p=\hbar k \]
\begin{align*}
\hat{H}\psi&=\frac{1}{2m}\hat{p}^{2}\psi\\
&=\kakimachigai{C\frac{\hbar^{2}k^{2}}{2m}\psi}{\frac{\hbar^{2}k^{2}}{2m}\psi}
\end{align*}
\[ E=\frac{\hbar^{2}k^{2}}{2m} \]
\toi{3}{○}{記述}
\[ \comm{\hat{H}}{\hat{p}}=\ii\hbar\cdot(-F)\neq0 \]
より同時観測不可能。\nuke{背理法}
\naze{交換子が $0$ でないことだけでは，同時固有関数が1つも無いとは言えない（交換しない行列にも共通の固有ベクトルが1本だけあることがある）。同時固有関数 $\phi$（$\hat{H}\phi=E\phi$，$\hat{p}\phi=p\phi$）があったとすると $\comm{\hat{H}}{\hat{p}}\phi=pE\phi-Ep\phi=0$ となるが，左辺は $-\ii\hbar F\phi$ で $F\neq0$ なので $\phi=0$ となり，固有関数であることに反する，まで書く。}
%% kizuki
\item 「交換子が $0$ でないから同時に定められない」で止めない。要点は交換子が $0$ でない「定数」であること：同時固有関数に作用させると $0$ になるはずだが，$0$ でない定数をかけて $0$ になる関数は $0$ しか無いので矛盾する。交換子が演算子（行列）のときは，作用させて $0$ になる特定の状態がありうる。
\item \mae{第5章 問題7・10-30} 前回も「成り立たない」ことを示しただけで理由の1文を書き落とした。結論を書く前に「何を仮定すると何と矛盾するか」を1行書く。
